\cotefile{162-5}
\scanpage{39}
\figure{\begin{tikzpicture}[x=1cm,y=1cm]\draw[line width=.45pt] (0,.9) -- (0,.18) .. controls (1.8,.1) and (2.3,.18) .. (3,.08) .. controls (5,-.12) and (6.2,.07) .. (7,.02) -- (7.15,.28) -- (9,.18);\end{tikzpicture}}

\margin{\unsure{!} Attention : la catégorie $\underline{\mathrm{Hom}}_{\mathrm{grab}}$, qui s'identitifie à la catégorie definie par un $R\mathrm{Hom}_{Z}(C_x,C'_x)$; ce point devant être élucidé!}

de même les (Cat)-groupes(ou groupes à valeurs dans (Cat)) comme étant
les groupes semisimpliciaux qui ``sont''des catégories i.e. qui \deleted{se défini\unsure{\deleted{d}s|d}-}
\deleted{s\unsure{e|o}nt} satsifont la condition d'exactitude familière, \deleted{i}permettant d'en déterminer
tous les composants à l'aide de $C_o$\deleted{$_{\unsure{u|a}}$} et $C_1$ . Appliquant le tapis de
Dold, dans le cas commutatif, on trouve que les (Cat)-groupes abéliens
forment une catégorie équivalente à celle des complexes de chaines de
groupes abéliens qui sont de longueur 1:
\[
0 \leftarrow C_o \leftarrow C_1 \leftarrow 0 \quad .
\]
Quand on explicite l\unsure{\deleted{e}a|\deleted{a}e} (Cat)-groupe abélien associé à un tel complexe de
chaines, on trouve ``ma'' construction chérie d\unsure{\deleted{a}e|\deleted{e}a} la ``catégorie de Picard''
associée au complexe: objets ceux de $C_o$ , flêches de $x$ dans $y$ les $z$ dans $C_1$
tels que $dz=y-x$ . On trouve donc à isomorphisme ppès tous les groupes
abéliens de (Cat) ainsi. Un \uline{complément important} est que toute catégorie
de Picard \inserted{\deleted{(resp. et}} \deleted{commutative)}est équivalente (au sens technique qu'on ne précise
pas ici) à \deleted{la catégorie} un vrai (Cat)-groupe \inserted{\deleted{(resp. (Cat)-groupe}} \deleted{commutatif)} .
\inserted{Quand on \unsure{simplifie} mod ``\unsure{Grab}'' équivalence, cela correspond à travailler dans la \uline{catégorie} \uline{dérivée} des complexes Seuls invariants sont $H_0\,H_1$ (\unsure{qui jouent|ou l'invariant} $\unsure{\Pi|\pi}_0\,\unsure{\Pi|\pi}_1$)}
\margin{\deleted{\unsure{Quillen dit} mais cela contredit les résultats de \unsure{Verdier} \unsure{la fin} \unsure{de} \illegible{1}}}
\margin{Mais \unsure{attention|alors} si on part d'une ``cat. de Picard commut.'', elle ne se réalise que par un (Cat)-g. \uuline{non} \uline{comm.}}

On notera qu'on a une \deleted{\unsure{interp}} classification beaucoup plus restreinte
que pour les H-espaces (resp. les H-espaces homotopiquement comm.) mod
équivalence d'homotopie, de façon précise entre les groupes resp. les grou-
pes abéliens de la catégorie homotopique. La raison en est qu\inserted{'}on est beau\deleted{co}coup
plus restrictif dans le contexte catégorique pour le sens d\unsure{\deleted{a}e|\deleted{e}a} l'associativité
(resp. et de la commutativité), puisqu'on ne demande pas seulement
l'existence d'isomorphismes d'associ\deleted{\unsure{i}}a\deleted{i}tivité\deleted{,} (qui suffirait pour avoir un
groupe dans \inserted{la catégorie homotopique} $(\mathrm{Cat})' \approxeq (\mathrm{Ssimpl})'$, modulo l'existence d'unités et d'inverse)
mais qu'on suppose \uline{donné\deleted{e}s} ces isomorphismes, de façon à satisfaire à cer-
taines compatibilités: c'est ce qui permet, me semble-t-il, de prouver qu'à
\margin{\deleted{\uuline{\unsure{pas vrai|pourquoi}}}} \margin{(} équivalence près, on peut supposer qu'il s'agit de vrais groupes dans
(Cat), \deleted{resp. de vraie commutativité.}

\margin{38}
